\chapter{Theoretical Background}
\label{ch:background}

\section{Intrusion Detection and Network Flows}

A network intrusion detection system (NIDS) watches network traffic and
produces alerts when it judges that traffic to be malicious. Two design
philosophies dominate the field. \textit{Signature-based} systems such as
Snort and Suricata match traffic against a curated list of known-bad
patterns; they are precise but unable to detect novel attacks.
\textit{Statistical} or \textit{learning-based} systems flag traffic that
deviates from a learned model of normal behaviour. Modern enterprise
deployments combine both. For a broader survey of machine-learning
approaches to intrusion detection see Buczak and Guven~\cite{Buczak2016}.

The unit of analysis in this thesis is the \textit{network flow}. A flow
is a sequence of packets that share a 5-tuple (source IP, source port,
destination IP, destination port, protocol) within a defined time window.
CICFlowMeter~\cite{Lashkari2017}, the de-facto extractor for CICIDS2017,
computes 84 statistics per flow: packet counts, byte counts, inter-arrival
times, packet lengths, TCP flag counts. After leakage filtering
(Section~\ref{sec:leakage}) 73 features remain. They fall into five
families: timing (Flow Duration, IAT statistics), packet-length, packet-
count, header metadata (TCP flag counts, header lengths), and bulk-rate.
Different attack families occupy distinctly different regions of this
73-dimensional space, which is what makes flow-based detection feasible.

\plainbox{A flow is best thought of as a single conversation between two
computers --- a phone call, rather than a single word. A whole
conversation tells you more than any one syllable, which is why
flow-level features are the natural unit of input for a network IDS.}

\section{Models}

Four supervised classifiers are compared. \textbf{Logistic regression} is
a linear model: $P(y = 1 \mid x) = \sigma(w^\top x + b)$, where $\sigma$
is the sigmoid. Fast, interpretable, strong baseline when the decision
boundary is roughly linear, and naturally well-calibrated.
\textbf{Random forest} is a bagging ensemble: many trees on bootstrap
samples of training data, predictions averaged. The averaging makes the
forest stable across runs and robust to noisy features.
\textbf{XGBoost} (eXtreme Gradient Boosting) builds trees one at a time,
each new tree fitting the residual errors of the previous ensemble. Often
beats random forests on tabular data, sometimes at the cost of
overfitting. \textbf{LightGBM} uses leaf-wise growth and histogram-based
splits; typically faster on large datasets, often slightly better
calibrated.

\plainbox{Logistic regression draws a single straight line through
feature space. Random forest asks a hundred different decision-tree
experts and takes a vote. XGBoost and LightGBM build a chain of small
experts where each one corrects the previous one's mistakes.}

\section{Evaluation Metrics}

Choosing the right metric matters more than choosing the right model when
data is imbalanced. With 85\% benign traffic, a trivial ``benign always''
classifier already gets 85\% accuracy. Macro F1 is therefore the headline
metric: the harmonic mean of precision and recall, averaged across all
classes weighting each class equally.

Two threshold-free ranking metrics complement macro F1: ROC-AUC (area
under the receiver operating characteristic) measures how well the model
ranks attacks above benign across all thresholds; PR-AUC focuses on the
precision-recall trade-off and is more informative when the positive
class is rare. The most operationally meaningful metric, however, is
recall at a fixed false positive rate. A SOC has a finite alert budget;
the question is ``how many real attacks does the model catch when false
alarms are capped at one in a thousand?''.

A binary classifier outputs probabilities and needs a threshold. The
canonical 0.5 threshold hides per-model differences and fails under
temporal shift. We use Youden's J statistic, $J = \mathrm{TPR} -
\mathrm{FPR}$, evaluated on validation. Youden's J is prevalence-invariant
--- it produces the same threshold whether the test set has 5\% or 50\%
attacks, which matters because the day-based splits have widely varying
attack rates.

\section{Adversarial Examples and Statistical Tests}

An \textit{adversarial example}~\cite{Goodfellow2015} is a small,
deliberately crafted perturbation that causes the model to misclassify
its input. In the network setting an attacker shifts inter-arrival times
and packet lengths within a small budget so that the attack flow is
predicted as benign. Threat models vary by what the attacker knows
(white-box vs black-box) and how much they can change ($L^\infty$ vs
$L^2$ budgets, perturbable feature subsets). Chapter~\ref{ch:adversarial}
specifies the threat model used in this thesis precisely.

Two statistical tests support the claims later. \textbf{McNemar's test}
compares two classifiers on the same test set by counting cases where the
first is wrong but the second is right ($b$) versus the reverse ($c$);
the statistic $\chi^2 = (\lvert b - c\rvert - 1)^2 / (b + c)$ follows
$\chi^2$ with one degree of freedom under the null of equal error rates.
\textbf{Bootstrap confidence intervals} resample the test set with
replacement $B$ times and report the 2.5th--97.5th percentile of each
metric, providing a sampling-noise floor for headline numbers.
